% Generated draft: requires mathematical and source review.
\documentclass[11pt]{article}
\usepackage[margin=25mm]{geometry}
\usepackage{fontspec}
\setmainfont{Latin Modern Roman}
\usepackage{amsmath,amssymb,mathtools,unicode-math}
\setmathfont{Latin Modern Math}
\usepackage{enumitem,needspace,xurl,hyperref}
\setlength{\parindent}{0pt}
\setlength{\parskip}{5pt}
\setlength{\emergencystretch}{3em}
\newcommand{\N}{\mathbb{N}}
\newcommand{\Z}{\mathbb{Z}}
\newcommand{\Q}{\mathbb{Q}}
\newcommand{\R}{\mathbb{R}}
\newcommand{\C}{\mathbb{C}}
\newcommand{\F}{\mathbb{F}}
\newcommand{\A}{\mathbb{A}}
\newcommand{\PP}{\mathbb{P}}
\newcommand{\E}{\mathbb{E}}
\newcommand{\eps}{\varepsilon}
\newcommand{\dd}{\,\mathrm{d}}
\begin{document}
\section*{UnsolvedMath ID: 700016}
\textbf{Source problem code:} catalogue entry
\subsection{Definitions and mathematical statement}
For a group \(G\), let
\[
\operatorname{Sub}(G)=\{H:H\leq G\},\qquad
\operatorname{NSub}(G)=\{H\leq G:(\forall g\in G)\ gHg^{-1}=H\}.
\]
The subgroup lattice is ordered by inclusion, with
\[
H\wedge K=H\cap K,\qquad H\vee K=\langle H\cup K\rangle.
\]
The normal subgroup lattice is the sublattice consisting of \(\operatorname{NSub}(G)\), with these inherited operations.

The question of Yu. L. Ershov (1965), Kourovka Notebook Problem 1.20, is to characterize the groups, or classes of groups, for which the lattice of normal subgroups is first-order definable in the lattice of all subgroups.

Under the reading that the normal subgroups must form a parameter-free definable subset of the subgroup lattice, a precise formulation for an individual group \(G\) is the existence of a first-order formula \(\varphi(x)\) in the language \(\{\subseteq\}\), with no parameters and with all variables ranging over subgroups, such that
\[
\forall H\in\operatorname{Sub}(G),\qquad
\bigl(\operatorname{Sub}(G),\subseteq\bigr)\models\varphi(H)
\quad\Longleftrightarrow\quad H\in\operatorname{NSub}(G).
\]
The order language is equivalent here to the lattice language, since the meet and join operations are definable from the order and the order is definable from the lattice operations.

For a class \(\mathcal C\) of groups, the uniform version asks for a single such formula satisfying
\[
\exists\varphi(x)\ \forall G\in\mathcal C\ \forall H\in\operatorname{Sub}(G),\qquad
\bigl(\operatorname{Sub}(G),\subseteq\bigr)\models\varphi(H)
\quad\Longleftrightarrow\quad H\trianglelefteq G,
\]
where \(\exists\varphi(x)\) is a metamathematical quantifier over first-order formulas. The task is to characterize the groups and classes satisfying the intended definability requirement.

The supplied wording does not specify whether parameters are permitted, whether one formula must work uniformly throughout a class, or whether definability means defining the actual subset of normal subgroups rather than interpreting an abstract copy of its lattice. Thus the displayed parameter-free, uniform formulation for classes is an explicit possible reading, and these qualifications require resolution before it can be treated as the definitive formulation. No restriction to finite groups or to any other particular class is supplied.
\subsection{Short English statement}
Determine which groups, and which classes of groups, allow their normal subgroups to be recognized by a first-order formula using only the structure of the subgroup lattice. The source leaves parameter use, uniformity across a class, and the precise meaning of lattice definability unspecified.
\subsection{Sources}
\begin{itemize}
\item Local UnsolvedMath registry, supplied complete source record: permanent ID \(700016\), exact source problem code \texttt{catalogue entry}, title \emph{Kourovka Notebook Problem 1.20}; upstream revision \texttt{d371ee917ac9647fbb06cd6fc529deec6f3dda82}. The supplied registry labels the problem \texttt{open}; this is a source label, not an independently verified assessment.
\item \emph{The Kourovka Notebook (Unsolved Problems in Group Theory)}, Problem 1.20, proposed by Yu. L. Ershov in the first issue (1965). The supplied background identifies the consulted edition as the 21st issue edition (2026), \texttt{arXiv:1401.0300v48}. The supplied collection description names E. I. Khukhro and V. D. Mazurov as editors. These bibliographic details and the source wording are supplied evidence and have not been independently checked.
\end{itemize}
\end{document}
